% GENERATED FILE. Edit canonical lessons, then run npm run generate:books.
% canonical-source-hash: fnv1a64:8313a8e496e72b1e
% canonical-lessons: HI-C172-rokna, HI-C172-bhagna, HI-C172-carhna, HI-C172-utarna, HI-C172-ghumna

\chapter{To stop (someone or something), To run away, To climb, To get down, To stroll}
\label{ch:hi-to-stop-someone-or-something-to-run-away-to-climb-to-get-down-to-stroll}
\hlchaptermodality{172}

\begin{chapteropening}
\textbf{By the end of this chapter:} \emph{I can say to stop (someone or something), to run away, to climb, to get down and to stroll.}

Complete the last lesson of chapter 172: I can say to stop (someone or something), to run away, to climb, to get down and to stroll.
\end{chapteropening}

\section[\texorpdfstring{\hi{रोकना} (roknā)}{रोकना (roknā)}]{\hi{रोकना} (roknā) — to stop (someone or something)}

\label{lesson:HI-C172-rokna}



\begin{quote}
Before the new one: say the Hindi for to hide, then the Hindi for to stop.
\end{quote}

\subsection*{You'll want to know: \hi{रोकना}}
\textbf{\hi{रोकना}} — \emph{roknā} — \textquotedblleft{}to stop (someone or something)\textquotedblright{}.

\begin{cousinweb}
\textbf{\hi{रोकना}} (\emph{roknā}) is a tadbhava, a word inherited and reshaped by sound change. It comes from \emph{*rokk}, \textquotedblleft{}stop\textquotedblright{}, a form CDIAL reconstructs from its descendants (CDIAL 10827). CDIAL calls its link with Sanskrit \emph{rudh-}, \textquotedblleft{}to obstruct\textquotedblright{}, not clear.
\end{cousinweb}

\begin{grammarlens}[title={What you've built}]
One more word for this chapter.
\end{grammarlens}

\subsection*{Your turn}
\begin{itemize}
\raggedright
  \item \emph{Say it:} \emph{roknā}
  \item \emph{Say it:} \emph{roknā}, once more
  \item \emph{Say it:} say \emph{roknā}
  \item \emph{Recall:} say the Hindi for to hide, then the Hindi for to stop, then say \emph{roknā} again
\end{itemize}

\begin{tcolorbox}[breakable,colback=teal!4,colframe=teal!35!black,title={Before you move on}]
What does \hi{रोकना} mean? (\textquotedblleft{}To stop (someone or something)\textquotedblright{}.) Say it once more.
\end{tcolorbox}
\section[\texorpdfstring{\hi{भागना} (bhāgnā)}{भागना (bhāgnā)}]{\hi{भागना} (bhāgnā) — to run away}

\label{lesson:HI-C172-bhagna}



\begin{quote}
Before the new one: say the Hindi for to stop, then the Hindi for to stop.
\end{quote}

\subsection*{You'll want to know: \hi{भागना}}
\textbf{\hi{भागना}} — \emph{bhāgnā} — \textquotedblleft{}to run away\textquotedblright{}.

\begin{cousinweb}
\textbf{\hi{भागना}} (\emph{bhāgnā}) is a tadbhava, a word inherited from Sanskrit and reshaped by sound change. It comes from \emph{bhajyate}, \textquotedblleft{}is broken\textquotedblright{} (CDIAL 9361). Bloch linked the sense \textquotedblleft{}flee\textquotedblright{} with a different old verb; CDIAL finds that unnecessary, since \textquotedblleft{}broken\textquotedblright{} could come to mean \textquotedblleft{}routed, fled\textquotedblright{}.
\end{cousinweb}

\begin{grammarlens}[title={What you've built}]
One more word for this chapter.
\end{grammarlens}

\subsection*{Your turn}
\begin{itemize}
\raggedright
  \item \emph{Say it:} \emph{bhāgnā}
  \item \emph{Say it:} \emph{bhāgnā}, once more
  \item \emph{Say it:} say \emph{bhāgnā}
  \item \emph{Recall:} say the Hindi for to stop, then the Hindi for to stop, then say \emph{bhāgnā} again
\end{itemize}

\begin{tcolorbox}[breakable,colback=teal!4,colframe=teal!35!black,title={Before you move on}]
What does \hi{भागना} mean? (\textquotedblleft{}To run away\textquotedblright{}.) Say it once more.
\end{tcolorbox}
\section[\texorpdfstring{\hi{चढ़ना} (caṛhnā)}{चढ़ना (caṛhnā)}]{\hi{चढ़ना} (caṛhnā) — to climb, to board}

\label{lesson:HI-C172-carhna}



\begin{quote}
Before the new one: say the Hindi for to stop, then the Hindi for to run away.
\end{quote}

\subsection*{You'll want to know: \hi{चढ़ना}}
\textbf{\hi{चढ़ना}} — \emph{caṛhnā} — \textquotedblleft{}to climb, to board\textquotedblright{}.

\begin{cousinweb}
\textbf{\hi{चढ़ना}} (\emph{caṛhnā}) is a tadbhava, a word inherited and reshaped by sound change. It comes from \emph{*caḍhati}, \textquotedblleft{}rises\textquotedblright{}, a form CDIAL reconstructs from its descendants (CDIAL 4578).
\end{cousinweb}

\begin{grammarlens}[title={What you've built}]
One more word for this chapter.
\end{grammarlens}

\subsection*{Your turn}
\begin{itemize}
\raggedright
  \item \emph{Say it:} \emph{caṛhnā}
  \item \emph{Say it:} \emph{caṛhnā}, once more
  \item \emph{Say it:} say \emph{caṛhnā}
  \item \emph{Recall:} say the Hindi for to stop, then the Hindi for to run away, then say \emph{caṛhnā} again
\end{itemize}

\begin{tcolorbox}[breakable,colback=teal!4,colframe=teal!35!black,title={Before you move on}]
What does \hi{चढ़ना} mean? (\textquotedblleft{}To climb, to board\textquotedblright{}.) Say it once more.
\end{tcolorbox}
\section[\texorpdfstring{\hi{उतरना} (utarnā)}{उतरना (utarnā)}]{\hi{उतरना} (utarnā) — to get down, to get off}

\label{lesson:HI-C172-utarna}



\begin{quote}
Before the new one: say the Hindi for to run away, then the Hindi for to climb.
\end{quote}

\subsection*{You'll want to know: \hi{उतरना}}
\textbf{\hi{उतरना}} — \emph{utarnā} — \textquotedblleft{}to get down, to get off\textquotedblright{}.

\begin{cousinweb}
\textbf{\hi{उतरना}} (\emph{utarnā}) is a tadbhava, a word inherited from Sanskrit and reshaped by sound change. It comes from \emph{uttarati}, \textquotedblleft{}passes over; descends\textquotedblright{} (CDIAL 1770). CDIAL traces the sense from \textquotedblleft{}crosses over\textquotedblright{} through \textquotedblleft{}lands\textquotedblright{} to \textquotedblleft{}gets down\textquotedblright{}.
\end{cousinweb}

\begin{grammarlens}[title={What you've built}]
One more word for this chapter.
\end{grammarlens}

\subsection*{Your turn}
\begin{itemize}
\raggedright
  \item \emph{Say it:} \emph{utarnā}
  \item \emph{Say it:} \emph{utarnā}, once more
  \item \emph{Say it:} say \emph{utarnā}
  \item \emph{Recall:} say the Hindi for to run away, then the Hindi for to climb, then say \emph{utarnā} again
\end{itemize}

\begin{tcolorbox}[breakable,colback=teal!4,colframe=teal!35!black,title={Before you move on}]
What does \hi{उतरना} mean? (\textquotedblleft{}To get down, to get off\textquotedblright{}.) Say it once more.
\end{tcolorbox}
\section[\texorpdfstring{\hi{घूमना} (ghūmnā)}{घूमना (ghūmnā)}]{\hi{घूमना} (ghūmnā) — to stroll, to wander}

\label{lesson:HI-C172-ghumna}



\begin{quote}
Before the new one: say the Hindi for to climb, then the Hindi for to get down.
\end{quote}

\subsection*{You'll want to know: \hi{घूमना}}
\textbf{\hi{घूमना}} — \emph{ghūmnā} — \textquotedblleft{}to stroll, to wander\textquotedblright{}.

\begin{cousinweb}
\textbf{\hi{घूमना}} (\emph{ghūmnā}) is a tadbhava, a word inherited and reshaped by sound change. It comes from \emph{*ghummati}, \textquotedblleft{}revolves\textquotedblright{}, a form CDIAL reconstructs from its descendants (CDIAL 4485).
\end{cousinweb}

\begin{grammarlens}[title={What you've built}]
One more word for this chapter.
\end{grammarlens}

\subsection*{Your turn}
\begin{itemize}
\raggedright
  \item \emph{Say it:} \emph{ghūmnā}
  \item \emph{Say it:} \emph{ghūmnā}, once more
  \item \emph{Say it:} say \emph{ghūmnā}
  \item \emph{Recall:} say the Hindi for to climb, then the Hindi for to get down, then say \emph{ghūmnā} again
\end{itemize}

\begin{tcolorbox}[breakable,colback=teal!4,colframe=teal!35!black,title={Before you move on}]
What does \hi{घूमना} mean? (\textquotedblleft{}To stroll, to wander\textquotedblright{}.) Say it once more.
\end{tcolorbox}
